\section{Calibration and the risk premium}\label{sec:calibration}

The estimation target is the joint behaviour of USDC during the March 2023 Silicon Valley Bank episode in hourly secondary-market data, supplemented by a baseline-regime window covering normal trading conditions in 2022. The model is the one-class representative-holder specification of Appendix~\ref{app:formal_model}. We estimate a seven-parameter vector by weighted matching of the six moments reported below; the free-parameter vector, the loss function, and the Nelder-Mead optimisation are documented in Appendix~\ref{app:calib_details}, which reports the converged vector $\phi^*$ of \eqref{eq:phi_star} at a final loss of $1.82$. The identification of the calibrated parameters is summarised in Section~\ref{sec:identification} and detailed in Appendix~\ref{app:calib_details}.

\subsection{Data}\label{sub:calibration_data}

The price data is a high-frequency (hourly close) series for the USDC-USD pair.\footnote{Source: CryptoCompare hourly close, accessed from the CryptoCompare public API.} The baseline-regime moments cover calendar year 2022; the stress-regime moments cover the 96-hour window from 22:00 UTC on 2023-03-10 (the SVB statement announcement) to 22:00 UTC on 2023-03-14 (Federal Reserve emergency lending facility announcement). We compute six moments per asset: (i) baseline mean close, (ii) baseline standard deviation of close, (iii) baseline frequency of close below 0.975 (the large-depeg threshold), (iv) stress regime minimum close, (v) stress regime mean absolute depeg in basis points, and (vi) stress frequency of close below 0.975. The empirical values are reported in Table~\ref{tab:moments_targets}.

\begin{table}[H]
\centering
\caption{Empirical target moments for the USDC calibration. Baseline moments are computed over calendar year 2022 (hourly close); stress moments over the 96-hour SVB window 2023-03-10 to 2023-03-14.}\label{tab:moments_targets}
\begin{tabular}{llr}
\toprule
Regime & Moment & Value \\
\midrule
Baseline & Mean close $\bar p_b$                                 & $1.00001$       \\
Baseline & Standard deviation of close $\hat\sigma(p_b)$          & $0.000754$      \\
Baseline & Frequency of close $< 0.975$, $\widehat{\Pr}(p_b<0.975)$ & $0.000$       \\
\midrule
Stress   & Minimum close $\min p_s$                              & $0.9022$        \\
Stress   & Mean absolute depeg $\overline{|p_s-1|}$               & $240.06$\,bps   \\
Stress   & Frequency of close $< 0.975$, $\widehat{\Pr}(p_s<0.975)$ & $0.396$       \\
\bottomrule
\end{tabular}
\end{table}

\subsection{Calibration result and interpretation}\label{sub:interp}

The calibration converges to a final loss of $L(\phi^*) = 1.82$ at the parameter vector $\phi^*$ reported in \eqref{eq:phi_star} of Appendix~\ref{app:calib_details}. Four of the six target moments are matched within reading tolerance; the two missed moments are the baseline standard deviation and the stress large-depeg frequency. The calibrated vector recovers the structural microstructure block ($\psi, \nu, \mu_q, \zeta$) and the signal noise $\sigma$ at values that match the observed price dynamics on the four matched moments, while the two misses reflect high-frequency exchange microstructure outside the redemption-coordination object (the residual diagnostics are in the internet appendix). The structural risk-premium decomposition reported in Section~\ref{sec:identification} is driven by the well-identified parameters and is robust to the residuals on the unfit moments.
